\section*{第 \ref{ch:b3:petrological-thermodynamics} 章——岩石热力学}

\begin{solution}{exo:b3:petrological-thermodynamics:1}
$\Delta_rH = -3029.3 - 910.7 + 3935.0 = \qty{-5.0}{kJ}$；$\Delta_rS = 133.5 +
41.5 - 207.4 = \qty{-32.4}{J.K^{-1}}$；$\Delta_rV = 60.40 + 22.69 - 100.07 =
\qty{-16.98}{cm^3} = \qty{-1.698}{J.bar^{-1}}$。斜率 $32.4/1.698 =
\qty{19.1}{bar/K}$。在 \qty{873}{K} 时，$P = 1 + (873.15\times(-32.4) +
5000)/(-1.698) = \qty{13.7}{kbar}$。
\end{solution}

\begin{solution}{exo:b3:petrological-thermodynamics:2}
$c = 2$ 且 $p = 3 > c$：不满足相律。三相只能在
\ce{Mg2SiO4 + SiO2 <=> 2 MgSiO3} 的单变线上共存；随机取得的岩石
并不位于该线上，因此该组合记录的是未完全进行的反应，例如
镁橄榄石被一圈顽火辉石包裹，使其无法与
石英接触。
\end{solution}

\begin{solution}{exo:b3:petrological-thermodynamics:3}
$a_{\text{钙铝榴石}} = 0.20^3 = 0.008$；$a_{\text{镁铝榴石}} = 0.15^3 =
\num{3.4e-3}$；$a_{\text{钙长石}} = 0.40$。
\end{solution}

\begin{solution}{exo:b3:petrological-thermodynamics:4}
$T_c = 15\,000/(2\times8.314) = \qty{902}{K}$；$W = 2\times8.314\times1000
= \qty{16.6}{kJ.mol^{-1}}$。
\end{solution}

\begin{solution}{exo:b3:petrological-thermodynamics:5}
$\Delta_rG = -5432.2 + 2\times(-2088.8) - 730.2 + 3\times3478.1 =
\qty{94.3}{kJ}$；$P = 1 + 94\,300/6.61 = \qty{14.3}{kbar}$。实验线
给出 $22.80\times726.85 - 1093 = \qty{15.5}{kbar}$，
比其高 \qty{1.2}{kbar}。
\end{solution}

\begin{solution}{exo:b3:petrological-thermodynamics:6}
$RT = \qty{7483}{J.mol^{-1}}$。当 $x = 0.5$ 时：$7483\ln0.5 + 20\,000/4 = -5187 +
5000 = \qty{-187}{J.mol^{-1}}$。当 $x = 0.111$ 时：$7483\times(0.111\ln0.111 +
0.889\ln0.889) + 20\,000\times0.111\times0.889 = -2609 + 1974 =
\qty{-635}{J.mol^{-1}}$，由对称性，0.889 时结果相同。
等量的两支整体成分为 0.5，吉布斯能为
\qty{-635}{J.mol^{-1}}：分离使其降低 \qty{448}{J.mol^{-1}}。
\end{solution}

\begin{solution}{exo:b3:petrological-thermodynamics:7}
$c + 2 = 5$ 个相，5 条线。(D) $A + C = 2B$：$0 + 2/3 = 2\times1/3$。(A)
$B + D = 2C$：$1/3 + 1 = 2\times2/3$。(C) $2A + D = 3B$：$0 + 1 = 3\times1/3$。
(B) $A + 2D = 3C$：$0 + 2 = 3\times2/3$。
\end{solution}

\begin{solution}{exo:b3:petrological-thermodynamics:8}
由杠杆规则，成分为 0.889 的相占 $(0.3 - 0.111)/(0.889 - 0.111) =
0.243$，即 \qty{24}{\percent}，成分为 0.111 的相占 \qty{76}{\percent}。
\end{solution}

\begin{solution}{exo:b3:petrological-thermodynamics:9}
$\partial P/\partial X_{\text{an}} = -3RT/(6.21X_{\text{an}}) =
-3\times8.314\times900/(6.21\times0.30) = \qty{-12000}{bar}$，即每单位
$X_{\text{an}}$：误差 0.01 时，压强变化 \qty{120}{bar}。
\end{solution}

\begin{solution}{exo:b3:petrological-thermodynamics:10}
$G'_{\text{mix}} = RT\ln[x/(1 - x)] + W(1 - 2x)$，且 $G''_{\text{mix}} =
RT/[x(1 - x)] - 2W$，在 $x(1 - x) = RT/(2W)$ 处为零。当 $x = 0.3$ 时：$T =
2\times20\,000\times0.21/8.314 = \qty{1010}{K}$；\omterm{def:b3:petrological-thermodynamics:solvus}{固溶线}在
$20\,000\times0.4/(8.314\ln(7/3)) = \qty{1136}{K}$ 处达到。两者之间，
晶体处于亚稳态：分离会降低其能量，但
微小涨落会使能量升高，因此第二相必须成核。
\end{solution}

\begin{solution}{exo:b3:petrological-thermodynamics:11}
由于 $G_{\text{mix}}(x) = G_{\text{mix}}(1 - x)$，公切线
与曲线相切于两个对称点 $x$ 和 $1 - x$，且为水平线：
$G'_{\text{mix}}(x) = 0$，即所求方程。解得 $T(x) = W(1 -
2x)/[R\ln((1 - x)/x)]$。令 $x = 1/2 - \varepsilon$，$1 - 2x = 2\varepsilon$
且 $\ln((1/2 + \varepsilon)/(1/2 - \varepsilon)) \approx 4\varepsilon$，故
$T \to W/(2R)$（当 $\varepsilon \to 0$ 时）；而对于 $x < 1/2$，$T(x)$ 随 $x$
趋近 1/2 而增大，因此该极限即为最大值。
\end{solution}

\begin{solution}{exo:b3:petrological-thermodynamics:12}
$(2.0 - 0.08\times10)/0.92 = \qty{1.30}{wt\percent}$。锰有利于
石榴子石生长：\omterm{def:b3:petrological-thermodynamics:pseudosection}{有效全岩成分}中锰含量降低后，边缘
生长所处的相区及等值线走向，不同于全岩分析
的预测。
\end{solution}

\begin{solution}{pb:b3:petrological-thermodynamics:1}
\textbf{1.} $\Delta_rH^\circ = -6640.0 + 2\times(-2593.8) - 910.7 + 3\times4234.0
= \qty{-36.3}{kJ}$；$\Delta_rS^\circ = 260.1 + 165.6 + 41.5 - 597.9 =
\qty{-130.7}{J.K^{-1}}$；$\Delta_rV = 125.28 + 88.30 + 22.69 - 302.37 =
\qty{-66.10}{cm^3} = \qty{-6.61}{J.bar^{-1}}$。

\textbf{2.} $-36\,300 + 298.15\times130.7 = \qty{+2.67}{kJ}$：为正，因此
钙长石在地表是稳定的。

\textbf{3.} $P = 1 + (900\times(-130.7) + 36\,300)/(-6.61) = \qty{12.3}{kbar}$；
斜率 $130.7/6.61 = \qty{19.8}{bar/K}$。

\textbf{4.} $1 + 81\,500/6.61 = \qty{12.3}{kbar}$ 在 \qty{900}{K} 时；$1 +
119\,800/6.61 = \qty{18.1}{kbar}$ 在 \qty{1200}{K} 时。

\textbf{5.} 在 626.85 和 \qty{926.85}{\celsius} 时，实验线给出
13.2 和 \qty{20.0}{kbar}：分别比表中数据高 0.9 和 \qty{1.9}{kbar}。

\textbf{6.} $1914\times6.61 = \qty{12.7}{kJ}$。综合不确定度为
$\sqrt{3.2^2 + (3\times2.0)^2 + (2\times2.0)^2 + 1.0^2} = \qty{8.0}{kJ}$：二者
之差为其 1.6 倍，对量热法而言是可能的，这也是该线
取自实验的原因。

\textbf{7.} $\Delta_rC_p$ 在固体中较小，其两项积分
部分相互抵消：表格使该线仅移动 26 和 \qty{112}{bar}。

\textbf{8.} \ce{SiO2}、\ce{Al2O3}、\ce{FeO}、\ce{MgO}、\ce{MnO}、\ce{CaO}、
\ce{Na2O}、\ce{K2O} 和 \ce{H2O}：$c = 9$。相：石榴子石、黑云母、
白云母、斜长石、蓝晶石、石英和流体，$p = 7$。

\textbf{9.} $F = 9 - 7 + 2 = 4$：压强、温度及两个成分变量。$p = 7 \le c = 9$：
相律成立，该矿物组合
可在压强–温度平面的一个区域内稳定存在。

\textbf{10.} 两个处于平衡的纯 \ce{Al2SiO5} 同质多象变体位于
单变线蓝晶石 = 夕线石上：这给出 $P$ 与 $T$ 之间的一个关系，
再由温度计即可求出压强。

\textbf{11.} 流体决定水的化学势：
自由组分 8 个，固相 6 个，$F = 8 - 6 + 2 = 4$，
且 $p \le c$ 仍成立。

\textbf{12.} $a_{\text{钙铝榴石}} = 0.10^3 = 0.001$，$a_{\text{钙长石}} =
0.30$；$\ln K = \ln(0.001/0.30^3) = -3.30$。

\textbf{13.} $K_D = (0.12/0.74)/(0.45/0.55) = 0.1622/0.8182 = 0.198$。

\textbf{14.} $T = (52\,635 + 2488)/(20.20 + 24.94\times1.618) = 55\,123/60.57
= \qty{910}{K}$，即 \qty{637}{\celsius}。

\textbf{15.} $P_0 = 22.80\times636.9 - 1093 = \qty{13429}{bar}$；$(3\times8.314
\times910.1/6.21)\ln(1/3) = \qty{-4016}{bar}$；$P = \qty{9.41}{kbar}$。

\textbf{16.} 在 \qty{9.41}{kbar}、\qty{644.2}{\celsius} 时，随后为
\qty{9.55}{kbar}；在该压强下为 \qty{644.9}{\celsius} 和
\qty{9.56}{kbar}，此后不再变化：\qty{645}{\celsius} 和
\qty{9.56}{kbar}。

\textbf{17.} \qty{641}{\celsius} 和 \qty{8.79}{kbar}。表格给出的端元线
约低 \qty{1}{kbar}（见第一部分）；
保留实验线。

\textbf{18.} 取 $X_{\text{an}} = 0.35$：\qty{642}{\celsius} 和
\qty{8.93}{kbar}，低 \qty{0.63}{kbar}。

\textbf{19.} 在 \qty{918}{K} 时：$P = 1 + (918\times12.6 - 7700)/0.571 =
\qty{6.8}{kbar}$，低于 \qty{9.56}{kbar}：蓝晶石，与观察一致。

\textbf{20.} $z = \num{9.56e8}/(2800\times9.81) = \qty{34.8}{km}$。

\textbf{21.} $\Delta P = (3\times8.314\times918/6.21)\ln(0.07/0.10) =
\qty{-1.3}{kbar}$：核部大约在较浅 \qty{4.8}{km} 处形成。石榴子石
在埋藏期间生长，其边缘记录了最深点。

\textbf{22.} 瓦尔登片岩达到的峰值压强为
\textbf{\qty[text-series-to-math=true]{9.6}{kbar}}，温度
\qty{645}{\celsius}，埋深约 \qty{35}{km}。
\end{solution}
